% Hide colored link boxes while keeping links clickable
\hypersetup{hidelinks}

\chapter{Certification Requirements}

\section{Significant Criteria}
This chapter describes major criteria across European countries for two popular ancillary services: \textit{FCR} and \textit{aFRR}. These parameters are either predefined before tests or calculated during tests to demonstrate the strength required to take actions in the power system.

The following table summarizes the significant criteria and indicates similarities and dissimilarities between countries.

\begin{longtable}{|p{3cm}|p{12cm}|}
    \caption{Major Criteria for Prequalification Test}\label{tab:criteria1}\\
    \hline
    \textbf{Criteria} & \textbf{Description} \\
    \hline
    \endfirsthead

    \hline
    \textbf{Criteria} & \textbf{Description} \\
    \hline
    \endhead

    \hline
    \endfoot

    \hline
    \endlastfoot

    Activation Speed & This parameter defines the time taken for the system to respond to a frequency signal determined by a frequency meter. It has three phases to activate the reserve: \textit{Initial Activation Time}, \textit{Half Activation Time} and \textit{Full Activation Time}. This time limit indicates the maximum time allowed for the response, and if any one of these three phases fails to meet the time limit, the test will be considered unsuccessful. \\
    \hline
    Capacity & This parameter explains the minimum and maximum capability of the system to generate power when the system lacks frequency, or to consume power when the system overgenerates and needs to reduce excess power. It is measured in MW. For the capacity market, this criterion indicates the potential allocated power reserved for future use as backup. \\
    \hline
    Delivery Duration & This parameter indicates how long the facility or system must continuously provide the ancillary service after activation. It is measured in hours or minutes and defines the minimum duration requirement for service delivery. \\
    \hline
    Non-LER & This is known as Non-Limited Energy Reservoir in the prequalification test, which determines the availability of the system to provide reserves continuously for a certain amount of time without any intervention. It is measured in hours, and interested participants need to clarify the longevity of their system, since an asset with limited energy requires special consideration and testing rules to be qualified for a market bid. \\
    \hline
    Meter Accuracy & Meter accuracy is the precision of the measurement of the frequency signal. It is measured as a percentage and indicates the maximum error allowed for the frequency meter to be qualified for a market bid. \\
    \hline
    Error Margin & This parameter has two values: \textit{Min Error Margin} and \textit{Max Error Margin}. These values indicate the acceptable range of response around the target value. For up-regulation, the maximum value is much more flexible than the minimum value, and vice versa for down-regulation. The TSO provides flexibility in terms of error margin to encourage interested bidders with their available resources. \\
    \hline
    Sampling Time & Sampling time indicates the gap between two consecutive measurements of the frequency signal. It is measured in seconds. This parameter is closely related to the frequency measurement; the higher the frequency, the shorter the sampling time required to capture one complete cycle of the signal. \\
    \hline
    Controller & The controller differs from other criteria since it is not a technical parameter but rather a requirement for the location of the asset. The controller works alongside the frequency meter to automatically send a response to the asset and trigger the activation of the reserve service. It receives a signal from the frequency meter as input, checks the frequency deviation, and determines an action based on the calculated decision within it. There are two types of controllers: \textit{Central Controller} and \textit{Local Controller}. A central asset is a cloud-based simulation software that receives a signal from the frequency meter remotely and sends a response to the asset for activation. A local controller means a controller with a frequency meter that can be located at the connection point of the transmission system. \\
    \hline
    Product Type & This criterion is mainly related to the regulation type or direction of response. There are some fixed categories for product type such as \textit{Symmetric} and \textit{Asymmetric Up Regulation} and \textit{Asymmetric Down Regulation}. The symmetric product type means that bidders must be able to provide both up and down regulation simultaneously as a sine wave signal. In contrast, the asymmetric product type is supplied by bidders for either up or down regulation. It brings significant differences to the rules of the prequalification test, since the symmetric product type must fulfil requirements for both directions together. \\
\end{longtable}


Though there is a set of test criteria enforced by the TSO, not all prerequisites apply to both FCR and aFRR markets. FCR has more requirements than aFRR. The following table shows what FCR and aFRR markets require in terms of criteria for the prequalification test. Some criteria need minimum and maximum values to set thresholds for specific parameters.

\begin{table}[htbp]
    \centering
    \begin{tabular}{|c|c|c|}
        \hline
        \textbf{Criteria} & \textbf{FCR} & \textbf{aFRR} \\
        \hline
        Activation Speed & Yes & Yes \\
        \hline
        Capacity & Yes & Yes \\
        \hline   
        Delivery Duration & Yes & Yes \\
        \hline  
        LER & Yes & Yes \\
        \hline
        Non-LER & Yes & Yes \\
        \hline
        Frequency Meter and Accuracy & Yes & N/A \\
        \hline
        Error Margin & Yes & Yes \\
        \hline
        Sampling Time & Yes & N/A \\
        \hline
        Central Asset & Yes & No \\
        \hline  
        Product Type & Yes & Yes \\
        \hline 
    \end{tabular}
    \caption{Significant Criteria for FCR and aFRR across European countries}
    \label{tab:criteria2}
\end{table}

This table indicates what information is required for each market to handle real-time dynamic cases while balancing frequency deviations.



\section{Importance of Certification}
Certification is the main focus of this study, providing validation and authorization for successful bidders. It refers to the mandatory prequalification tests and approval process that a reserve unit must undergo before it can participate in the market. Strategically, this process means:
\begin{itemize}
    \item \textbf{Mandatory Technical Verification:} A reserve unit must demonstrate through prequalification tests that its control behaviour and performance meet specific technical requirements, such as activation speed, stability and endurance. Controlling a large system, where errors can lead to both financial loss and physical damage, is a complex task.
    \item \textbf{Prerequisite for Market Entry:} A Balancing Service Provider (BSP) can only offer capacity to the hourly or annual markets once the TSO has formally approved the capacity based on the test results.
    \item \textbf{Risk Mitigation:} Certification helps mitigate risks associated with non-performance or underperformance of reserve units, which can lead to penalties or financial losses for both the provider and the system operator.
    \item \textbf{Encouraging Investment:} A clear certification process can encourage investment in reserve capacity by providing confidence to investors that their assets will be recognised and compensated for their contributions to grid stability.
\end{itemize}




\begin{landscape}

    \section{Comparative Analysis of Certification Process}
The following discussion covers the certification process across three major zones in Europe. These zones are Northern Europe, the Continental Europe Synchronous Area, and Southern Europe. Initial similarities are visible in the types of ancillary services offered. In general, the Nordic region employs more frequency control facilities than central and southern Europe. A brief description and comparison across the three zones are given below:
    \subsection{Zone 1: Nordic Synchronous Area}
    All Nordic countries such as Finland, Sweden, Norway and the eastern part of Denmark actively engage in the electricity market. This area is highly regulated by renewable energy, which makes production levels unpredictable and causes rapid changes in frequency. As a result, this synchronous area has introduced more ancillary services and extended the market to new bidders. The popular ancillary services are FFR, FCR-N, FCR-D up-regulation, FCR-D down-regulation, aFRR, and mFRR.
    
    In this section, the prequalification test requirements are included only for Eastern Denmark, also known as DK2. The permission to provide ancillary services has a constraint in terms of participants' roles for Denmark. A BSP can only access the ancillary service market with a limited energy supply, which covers FFR and FCR-D. A table of technical requirements for the FCR-D prequalification test in Eastern Denmark is provided below.


\tiny
\setlength{\tabcolsep}{2pt}
\renewcommand{\arraystretch}{1.3}
\begin{longtable}{|p{1.5cm}|p{2.5cm}|p{2.5cm}|p{1.8cm}|p{2.5cm}|p{1.5cm}|p{4cm}|p{2.5cm}|p{2.5cm}|p{2cm}|}
    
\caption{Zone 1: Nordic Synchronous Area --- comparison of prequalification criteria}\label{tab:Nordic}\\

\hline
\textbf{Country} & \textbf{Activation Speed} & \textbf{Capacity} & \textbf{Delivery Duration} & \textbf{Non-LER Duration} & \textbf{Meter Accuracy} & \textbf{Error Margin} & \textbf{Frequency Meter} & \textbf{Sampling Time} & \textbf{Product Type} \\
\hline
\endfirsthead

\hline
\textbf{Country} & \textbf{Activation Speed} & \textbf{Capacity} & \textbf{Delivery Duration} & \textbf{Non-LER} & \textbf{Meter Accuracy} & \textbf{Error Margin} & \textbf{Frequency Meter} & \textbf{Sampling Time} & \textbf{Product Type} \\
\hline
\endhead

Eastern Denmark & 
\textbf{5\% act.:} 2.5s \newline
\textbf{86\% act.:} 7.5s \newline & 
\textbf{Min:} 0.1 MW~\cite{energinet-market-requirements-dk2-fcrd} \newline
\textbf{Max:} N/A & 
20 min~\cite{energinet-market-requirements-dk2-fcrd} & 
\textbf{Asym. (Up / Down):} 2h \newline
\textbf{Asym. (Up \& Down):} 4h & 
\textbf{Acc.:} > 10 mHz & 
\textbf{Fast Ramp Test (Dynamic FCR-D): }
\newline
1. Under-delivery limit: 5\% 
\newline
2. Over-delivery limit: 20\%

\vspace{0.2cm}
\textbf{Ramp Static Test (Static FCR-D): }
\newline
1. Under-delivery limit: 5\% 
\newline
2. Over-delivery limit: 10\%

\vspace{0.2cm}
\textbf{Linearity Test (Steady State Response): }
\newline
Upward Regulation:  
\newline 1. Under-delivery limit: -5\% 
\newline 2. Over-delivery limit: +10\% 
\vspace{0.09cm} \newline
Downward Regulation: 
\newline 1. Under-delivery limit: -10\%
\newline 2. Over-delivery limit: +5\%
&
Both local and central frequency meters are acceptable, though a local frequency meter is recommended for each unit to ensure higher security of delivery. For Denmark, a central meter means that multiple units can share the same frequency meter to create redundancy in case of power outages or equipment failure.
& 
\textbf{Frequency: }1 per second
\textbf{Power: }1 per second
& \begin{minipage}[t]{\linewidth}
\begin{itemize}[leftmargin=*, nosep]
  \item Asymmetric Up
  \item Asymmetric Down
\end{itemize}
\vspace{-0.5\baselineskip}\cite{energinet-market-requirements-dk2-fcrd}
\end{minipage} \\
\hline
\end{longtable}
\normalsize
\end{landscape}




\clearpage
\begin{landscape}


    \subsection{Zone 2: Continental Europe Synchronous Area}
In this section, some central European countries, namely Belgium, Austria, Netherlands, Switzerland, Czechia, Germany, and the Baltics, are covered. These countries have a remarkable business environment, well-established working processes, and opportunities for international traders, which foster great collaboration among diverse market actors. Continental Europe has more control over electricity production and frequency stabilisation because the major portion of generation comes from nuclear energy and fossil fuels. This indicates a high-inertia system where imbalances in supply and demand can be mitigated without advanced control systems. As a result, continental Europe maintains frequency balance with fewer ancillary services and a more moderate market policy than the northern part of Europe.

The popular ancillary services traded in the electricity market sector are FCR, aFRR, and mFRR. A data table on the FCR market in continental Europe is added below:

\tiny
\setlength{\tabcolsep}{2pt}
\renewcommand{\arraystretch}{1.3}
\setlength{\LTleft}{0pt}
\setlength{\LTright}{0pt}
\begin{longtable}{|p{1.8cm}|p{2.2cm}|p{2.2cm}|p{2.2cm}|p{2.5cm}|p{1.5cm}|p{2.2cm}|p{4cm}|p{2.5cm}|p{2cm}|}
\caption{Zone 2: Continental Europe Synchronous Area --- comparison of FCR prequalification criteria}\label{tab:continental_europe_fcr}\\
\hline
\textbf{Country} & \textbf{Activation Speed} & \textbf{Capacity} & \textbf{Delivery Duration} & \textbf{Non-LER Duration} & \textbf{Meter Accuracy} & \textbf{Error Margin} & \textbf{Frequency Meter} & \textbf{Sampling Time} & \textbf{Product Type} \\
\hline
\endfirsthead

\hline
\textbf{Country} & \textbf{Activation Speed} & \textbf{Capacity} & \textbf{Delivery Duration} & \textbf{Non-LER Duration} & \textbf{Meter Accuracy} & \textbf{Error Margin} & \textbf{Frequency Meter} & \textbf{Sampling Time} & \textbf{Product Type} \\
\hline
\endhead

Belgium & 
\textbf{Initial Act.:} $\leq$ 2 sec \newline
\textbf{Half Act.:} $\leq$ 15 sec \newline
\textbf{Full Act.:} $\leq$ 30 sec \newline
& 
\textbf{Min: }1 MW 
\newline \textbf{Max: }N/A
& To be inserted & 2h & 
\textbf{Acc.: }> 10mHz
& 
\textbf{Availability Test: }> 15\%
&
1. Local frequency measurements are mandatory. That means each Delivery Point (i.e., each site participating in the Frequency Containment Reserve or FCR service) must have its own frequency measurement equipment installed on-site.  
\newline \textbf{Exception: Virtual Delivery Points}
\newline 2.1 For Virtual Delivery Points (which are aggregations of smaller assets treated as one point): 
\newline 2.2 Only one frequency meter is required per Virtual Delivery Point. 
\newline 2.3 However, if a Balancing Service Provider (BSP) operates multiple Virtual Delivery Points, then: 
\newline \hspace{1em} 2.3.1 Each Virtual Delivery Point must have its own separate frequency meter. 
\newline \hspace{1em} 2.3.2 This is for redundancy, meaning to avoid single points of failure across multiple assets. 

\textbf{Administrative Note:} If a BSP wants a Providing Group to be treated as a Virtual Delivery Point, they must notify ELIA before any prequalification tests are conducted.
& N/A & 
\begin{itemize}[leftmargin=*, nosep]
  \item Symmetric 100mHz
  \item Symmetric 200mHz
  \item Asymmetric Up
  \item Asymmetric Down
\end{itemize}\\
\hline
Austria & 
\textbf{Half Act.:} $\leq 15\,\mathrm{s}$

\textbf{Full Act.:} $\leq 30\,\mathrm{s}$

& 
\textbf{Min: }$\pm 1$ MW
\textbf{Max: }Depends on prequalified capacity
\newline \textbf{Max indivisible: }$\pm 25$ MW
\textbf{Resolution: }1 MW
& 
\textbf{Min: }30 min
& 
To be inserted 
&
\textbf{Error} $\leq$ $\pm 5$ mHz 
& 
To be inserted 
& 
To be inserted 
& 2 sec & To be inserted \\

\hline
Netherlands 
& 
\textbf{Initial Act.:} $\leq$ 2 sec \newline
\textbf{Half Act.:} $\leq$ 15 sec \newline
\textbf{Full Act.:} $\leq$ 30 sec \newline
& 
\textbf{Min: }1 MW \newline
\textbf{Max: }150 MW \newline
\textbf{Unit size: }$\geq$ 0.1 MW \newline
\textbf{Inc. size: }0.1 MW \newline

All units in a group must have similar power.
\newline Maximum size for indivisible bids = 25 MW.
& To be inserted 
& 
\textbf{Asym. (Up/Down): } 2h
\newline
LER + non-LER = 4h (full group considered as non-LER)
& 
\textbf{Acc.: }$\geq$ 10mHz 
& 
To be inserted
&
Each asset requires a local frequency meter; a central controller may send signals only where no local intervention is needed.
& To be inserted & Symmetric \\
\hline
Switzerland 
& 
\textbf{Initial Act.:} $\leq$ 2 sec \newline
\textbf{Half Act.:} $\leq$ 15 sec \newline
\textbf{Full Act.:} $\leq$ 30 sec \newline
& 
\textbf{Min: }1 MW
& 
To be inserted 
& 
To be inserted & 
\textbf{Acc.: }> $\pm 10$mHz  
& 
To be inserted 
& 
There are 3 approaches for frequency measurement:
\newline 1. Local frequency measurements at least per network connection point (if possible, per technical unit).
\newline 2. Central controller with decentralised frequency measurements per connection point (assumes separate control).
\newline 3. Central controller with integrated frequency measurements of all technical units (assumes integrated control); the TSO selects the best result from option 2 or option 3.
& 
To be inserted &
Up/Down \\
\hline
Czechia 
& 
\textbf{Initial Act.:} $\leq$ 2 sec \newline
\textbf{Half Act.:} $\leq$ 15 sec \newline
\textbf{Full Act.:} $\leq$ 30 sec
& 
\textbf{Min: }1 MW
\newline \textbf{Max: }25 MW
& To be inserted
&
2h
& 
\textbf{Acc.: }$\geq$ 10mHz  
\vspace{0.2cm}
For backup frequency meter,
\newline \textbf{Acc.: }$\geq$ 50mHz  
& To be inserted 
& 
1. Units and aggregation blocks (ABs) used for backup must be equipped with frequency measurement at each connection point.
2. If one AB contains more than one connection point, centralised measurement may be used as the main measurement.
3. If condition 2 applies, then each connection point within the AB where the energy device is located must also be equipped with a backup measurement.
& To be inserted & 
Up/Down \\
\hline
Germany 
& 
\textbf{Initial Act.:} $\leq$ 2 sec \newline
\textbf{Full Act.:} $\leq$ 30 sec 
& To be inserted & To be inserted & 2h & \textbf{Acc.: }10mHz & 
10\% to 30\% depending on the state in which the activation occurs.
& 
One meter per grid connection point.
& 
4 sec 
& 
Symmetric $\pm 200$mHz \\
\hline
Baltics 
&
\textbf{Initial Act.:} $\leq$ 2 sec \newline
\textbf{Full Act.:} $\leq$ 30 sec
&
\textbf{Min: }1 MW
& To be inserted 
& 
2h & 
\textbf{Acc.: } $\pm 10$mHz
& 
$\pm 10$\%
& 
Local frequency meter mandatory.
& 
100 ms
& Symmetric $\pm 200$mHz \\
\hline
\end{longtable}
\normalsize
\end{landscape}




\clearpage
\begin{landscape}

    Another widely used market product in continental Europe is aFRR, known as automatic frequency restoration reserve. Most countries share notable similarities from the perspective of technical requirements. A data table comparing the aFRR market requirements in continental Europe is provided below:


\tiny
\setlength{\tabcolsep}{4pt}
\renewcommand{\arraystretch}{1.3}
\setlength{\LTleft}{0pt}
\setlength{\LTright}{0pt}
\begin{longtable}{|p{2.5cm}|p{2.5cm}|p{2.5cm}|p{2.5cm}|p{2.5cm}|p{4cm}|p{2.5cm}|p{2.5cm}|}
\caption{Zone 2: Continental Europe Synchronous Area --- comparison of aFRR prequalification criteria}\label{tab:continental_europe_afrr}\\
\hline
\textbf{Country} & \textbf{Activation Speed} & \textbf{Capacity} & \textbf{Delivery Duration} & \textbf{Non-LER Duration} & \textbf{Error Margin} & \textbf{Sampling Time} & \textbf{Product Type} \\
\hline
\endfirsthead

\hline
\textbf{Country} & \textbf{Activation Speed} & \textbf{Capacity} & \textbf{Delivery Duration} & \textbf{Non-LER Duration} & \textbf{Error Margin} & \textbf{Sampling Time} & \textbf{Product Type} \\
\hline
\endhead

Belgium & 
\textbf{Initial Act.:} $\leq$ 2 sec \newline
\textbf{Half Act.:} $\leq$ 15 sec \newline
\textbf{Full Act.:} $\leq$ 30 sec \newline
& 
\textbf{Min: }1 MW 
\newline \textbf{Max: }9999 MW
& To be inserted & 
\textbf{Asym. (Up/Down): }4h 
\newline SoC > 50\% at start
& 
2 MW threshold over 10 seconds serves as the error margin used for compliance or deviation tracking.
& 1 sec & 
\begin{itemize}[leftmargin=*, nosep]
  \item Asymmetric Up
  \item Asymmetric Down
\end{itemize}\\
\hline

Austria & 
\textbf{Reaction time:} 30 sec \newline
\textbf{Full Act.:} 5 min \newline
& 
\textbf{Min: }$\pm 1$ MW
\newline \textbf{Inc. size: }1 MW 
& 
To be inserted
& 
To be inserted 
& 
1. Under-delivery limit: 5\%
\newline 2. Over-delivery limit: 10\%
& 2 sec & To be inserted \\

\hline
Netherlands 
& 
\textbf{Initial Act.: }30 sec
& 
\textbf{Min: }1 MW
& 
To be inserted
& 
N/A
& 
1. Under-delivery limit: 10\%
\newline 2. Over-delivery limit: 20\%
& 10 sec 
& \begin{itemize}[leftmargin=*, nosep]
  \item Asymmetric Up
  \item Asymmetric Down
\end{itemize}\\

\hline
Switzerland 
& 
To be inserted
& 
\textbf{Min: }$\pm 5$ MW
\newline \textbf{Max: }$\pm 100$ MW
& 
To be inserted 
& 
To be inserted & 
2 conditions on error margin under/over delivery limit:

1. limit $\leq$ $\pm 5$\%, acceptable.

2. limit > $\pm 5$\%, the sum of all responses exceeding 5\% must be $\leq$ 1\% of total target power.
& 
1. 2 sec (recommended)
\newline 
2. 10 sec (acceptable) &
\begin{itemize}[leftmargin=*, nosep]
    \item Symmetric
    \item Asymmetric Up
    \item Asymmetric Down
\end{itemize} 
\\

\hline
Germany 
& 
\textbf{Initial Act.:} 30 sec \newline
\textbf{Full Act.:} 5 min 
& 
\textbf{Min: }1 MW
& To be inserted 
& 4h & 
5--10\% depending on the location of activation.
& 
4 sec 
& 
Up / Down \\


\hline
Baltics 
&
\textbf{Initial Act.:} 30 sec \newline
\textbf{Full Act.:} 5 min 
&
\textbf{Min: }1 MW
& To be inserted 
& 
2h & 
$\pm 10$\%
& 
10 seconds on asset level
& Up / Down  \\
\hline
\end{longtable}
\normalsize
\end{landscape}



    
%Zone 3



\clearpage
\begin{landscape}


    \subsection{Zone 3: Southern Europe Synchronous Area}

In this section, countries in the southern part of Europe and neighbouring regions are included. There are very few TSOs in this zone. Similar to continental Europe, the major portion of generation comes from nuclear energy and fossil fuels, indicating a high-inertia system where imbalances in supply and demand can be mitigated without advanced control systems. As a result, southern Europe maintains frequency stability with fewer ancillary services and a more moderate market policy than the northern part of Europe.

The popular ancillary services traded in the electricity market sector are FCR, aFRR, and mFRR. A data table on the FCR market in southern Europe is added below:

\tiny
\setlength{\tabcolsep}{2pt}
\renewcommand{\arraystretch}{1.3}
\setlength{\LTleft}{0pt}
\setlength{\LTright}{0pt}
\begin{longtable}{|p{1.8cm}|p{2.2cm}|p{2.2cm}|p{2.2cm}|p{2.5cm}|p{1.5cm}|p{2.2cm}|p{4cm}|p{2.5cm}|p{2cm}|}
\caption{Zone 3: Southern Europe Synchronous Area --- comparison of FCR prequalification criteria}\label{tab:south_europe_fcr}\\
\hline
\textbf{Country} & \textbf{Activation Speed} & \textbf{Capacity} & \textbf{Delivery Duration} & \textbf{Non-LER Duration} & \textbf{Meter Accuracy} & \textbf{Error Margin} & \textbf{Frequency Meter} & \textbf{Sampling Time} & \textbf{Product Type} \\
\hline
\endfirsthead

\hline
\textbf{Country} & \textbf{Activation Speed} & \textbf{Capacity} & \textbf{Delivery Duration} & \textbf{Non-LER Duration} & \textbf{Meter Accuracy} & \textbf{Error Margin} & \textbf{Frequency Meter} & \textbf{Sampling Time} & \textbf{Product Type} \\
\hline
\endhead

France & 
\textbf{Initial Act.:} $\leq$ 2 sec \newline
\textbf{Full Act.:} $\leq$ 30 sec \newline
& 
\textbf{Min: }1 MW 
\newline \textbf{Max: }25 MW
& 15 min & 2h & 
\textbf{Acc.: }> 10mHz
& 
To be inserted
&
Both local and centralised frequency measurement are permitted under certain conditions:

Local frequency measurement:
  1. Any site with FCR capacity > 1 MW.
  2. If a site reacts directly and proportionally to frequency deviation (for aggregators).

Centralised frequency measurement:
  1. At least one site is not activated using its own local frequency measurement (for aggregators).

Sites connected to the distribution grid (PDS):
  1. If sites do not have local frequency measurement:
        Once the total PFC capacity in a region exceeds 1.5 MW,
        at least 3 frequency measurements are required,
        located in 3 different departments,
        and measurement quality must be equivalent to control measurement.
& N/A & 
\begin{itemize}[leftmargin=*, nosep]
  \item Symmetric
  \item Asymmetric
\end{itemize}\\
\hline
\end{longtable}
\normalsize
\end{landscape}




\clearpage
\begin{landscape}

    Another widely used market product in southern Europe is aFRR, known as automatic frequency restoration reserve. Most countries share notable similarities from the perspective of technical requirements. A data table comparing the aFRR market requirements in southern Europe is provided below:


\tiny
\setlength{\tabcolsep}{4pt}
\renewcommand{\arraystretch}{1.3}
\setlength{\LTleft}{0pt}
\setlength{\LTright}{0pt}
\begin{longtable}{|p{2.5cm}|p{2.5cm}|p{2.5cm}|p{2.5cm}|p{2.5cm}|p{4cm}|p{2.5cm}|p{2.5cm}|}
\caption{Zone 3: Southern Europe Synchronous Area --- comparison of aFRR prequalification criteria}\label{tab:south_europe_afrr}\\
\hline
\textbf{Country} & \textbf{Activation Speed} & \textbf{Capacity} & \textbf{Delivery Duration} & \textbf{Non-LER Duration} & \textbf{Error Margin} & \textbf{Sampling Time} & \textbf{Product Type} \\
\hline
\endfirsthead

\hline
\textbf{Country} & \textbf{Activation Speed} & \textbf{Capacity} & \textbf{Delivery Duration} & \textbf{Non-LER Duration} & \textbf{Error Margin} & \textbf{Sampling Time} & \textbf{Product Type} \\
\hline
\endhead

France & 
\textbf{Initial Act.: }60 sec \newline
\textbf{Full Act.: }5 min
& 
\textbf{Min: }1 MW 
\newline \textbf{Max: }N/A
& To be inserted & 
To be inserted
& 
To be inserted
&  & 
\begin{itemize}[leftmargin=*, nosep]
  \item Symmetric
  \item Asymmetric Up
  \item Asymmetric Down
  \item Symmetric and Asymmetric
\end{itemize}\\
\hline

Spain & 
\textbf{Init. time:} 20 sec \newline
\textbf{Full Act.:} 15 min \newline
& 
& 
& 
& 
& 2 sec & To be inserted \\
\hline
\end{longtable}
\normalsize
\end{landscape}
